\documentclass[11pt,english]{article}
\usepackage[utf8]{inputenc}
\usepackage[a4paper,left=0.5in,right=0.5in,top=1in,bottom=1in]{geometry}
\usepackage{newpxtext}
\usepackage{graphicx}
\usepackage{caption}
\usepackage{subcaption}
\usepackage{float}
\usepackage{multicol}
\usepackage{fancyhdr}

\begin{document}
\begin{center}
\Large\textbf{An agent based model of Brenner´s theory of transition}
\end{center}

\begin{abstract}
    The results of the Brenner Debate produced an enhanced theory of transition from feudalism to capitalism based in rural $17^{th}$ Century England. A key tenant of this school of thought is the role of ``improvement" as both a motor in the divorce between peasants and their means of subsistence, and a consequence of newly created market imperatives both landlords and tenants were now subject too. To test the role of improvement ideology in the transition an agent based model is constructed. This incorporates tenants, landlords, ecological deterioration and rental markets. The critical role of rental market pressure is modelled through a transition from customary rent to leasehold rents. 
\end{abstract}
\begin{multicols}{2}
\normalsize
\section{Introduction}
According to \cite[p.10]{comninel2000english} feudalism may be defined as a mode of production in which the ruling class extracts surplus value from direct producers via non-economic means, often political and under the threat of force. In contrast, capitalism extracts surplus value from labourers, who are not direct producers and do not posses the means of production, via economic means. Capitalists must re-invest the extracted surplus value, in order to reproduce their class position, to keep up with increasing productivity. Through an analysis of the transition from one class structure to another, one might hope to gain insights into common identifiers which are associated with the collapse of a mode of production. Improvement ideology emerged in the mid $16^{th}$ Century composed of both practical methods to increase the output of a given piece of land, and as a method through which to better the condition of the English peoples, as argued by Captain Walter Blith.

\section{Marx on the transition and primitive accumulation}
\label{Marx on the transition}
Marx argued against the inevitability of capitalism, comparing it to the ``original sin'' in theology \cite[p.507]{marx2015capital}. Where capitalism's ``origin is supposed to be explained when it is told as an anecdote of the past" \cite[p.507]{marx2015capital}. Marx instead produced a theory based off the end result and its defining characteristics of wage co-modification and the separation of the means of production from wage-employees. For Marx the transformation of goods into capital requires a social transformation in the relations between the producer and the ruling class. This process Marx labels Primitive Accumulation. He identified the "expropriation of the self-supporting peasants" as a factor but also ''the destruction of rural domestic industry" \cite[p.531]{marx2015capital}. The former of which is the first step in the formation of capital, drawing a distinction between feudal wealth and capital. Marx states ''what enables money-wealth to become capital is the encounter, on one side, with free workers; and on the other side, with the necessaries and materials etc, which previously were in one way or another the property of the masses who have now become object-less, and are also free and purchasable" \cite{marx2005grundrisse}. Marx identifies this process of dispossession as having two components, firstly the ``forcible driving of the peasantry from [their] land" and secondly a ``usurpation of the common lands" \cite[p.514]{marx2015capital}. With both factors Marx uses Francis Bacon's ``History of the Reign of King Henry VII" to pinpoint historical legislation off the back of which much of this dispossession was carried out. For Marx these tendencies contributed to a final state of a ``degraded and almost servile condition of the mass of the people, the transformation of them into mercenaries, and of their means of labour into capital" \cite[p.511]{marx2015capital}. This degradation of the peasantry can be seen in the collapse of the independent yeoman who was replaced by tenants who depended on yearly leases. Producing ``a servile rabble dependent on the pleasure of the landlords" \cite[p.511]{marx2015capital}. This available labour was now ``transformed into material elements of variable capital" \cite[p.530]{marx2015capital} thereby introducing the final components for the transition, `` the peasant, expropriated and cast adrift, must buy their value in the form of wages, from his new master, the industrial capitalist" \cite[p.530]{marx2015capital}. For Marx it is a combination of social relations and material conditions that produces the banality of capitalist economic logic resulting in ``the subjugation of the labourer to the capitalist" \cite[p.523]{marx2015capital}. In his analysis Marx had a very limited mention of improvement language. Suggesting the ''revolution in the conditions of landed property was accompanied by improved methods of culture, greater co-operation, concentration of the means of production" \cite[p.530]{marx2015capital}. This lack of involvement of improvement within Marx's theory of transition could in part be explained by a lack of access to many of the primary sources used by the likes of Brenner and Woods in constructing their own theories. But also speaks to the fact that Marx was first and foremost a philosopher not a historian.

\section{Brenner's theory of transition}\label{Background and the Brenner debate}
With the publication of "Agrarian Class Structure and Economic Development in Pre-Industrial Europe" in 1976 Robert Brenner produced a crucial advance in the causes of the very first case study of the transition from feudalism to capitalism. It is in this very word "transition" that Brenner argued even Marxist historians such as Paul Sweezy or Maurice Dobb had been led astray \cite[p.41-53]{brenner1977origins}, emphasising that they were assuming an a priori existence of an embryonic capitalism within feudal society.Several theories of transition may be grouped as commercial models. The keystone of these theories is that capitalism arose as the barriers for its development were removed \cite[p.12]{wood2002origin}. Crucially, the driving factor behind this development was man's "propensity to truck, barter, and exchange" \cite{smith1963inquiry}. This force led to the division of labour and ultimately industrial capitalism. Brenner and Wood's fundamental critique of this approach is that it assumes a nascent agricultural or industrial capitalism within other modes of production, such that there is no real transition and capitalism was inevitable \cite[p.51]{wood2002origin}. Brenner's theory sought internal conflicts within feudalism, as opposed to an external impetus, which might produce its collapse. Wood offers an analysis that `` it [was] a matter of lords and peasants, in certain specific conditions peculiar to England, involuntarily setting in train a capitalist dynamic while acting in class conflict with each other to reproduce themselves as they were” \cite[p.52]{wood2002origin}. This focus on the class conflict nature of the transition is what distinguished Brenner and Wood's theory from other contemporary Political Marxists. Brenner identifies the internal logic of feudalism that inhibits surplus reinvestment. For the feudal lord it was often easier to find new avenues to "squeeze" \cite[p.74]{brenner1976agrarian} out more absolute surplus value from their subjects instead of increasing their labour efficiency. Due to the fact that a lord's power was directly correlated to size of the population under his control this lead to a focus on labour immobility, and maximising the ease with which he could extract surplus value from his peasants by non-economic means. Resulting in restrictions on land mobility and "unfree peasants [unable] to convey their land to other peasants without the lord's permission" \cite[p.50]{brenner1976agrarian}. The centuries previous to the $17^{th}$ Century were characterised by the end of serfdom in England, resulting in peasants whose lords extracted surplus value partly through the rent of their land. These were customary rents that decreased in value for the lord over time due to inflation. This loss in income revenue was counteracted through the imposition of arbitrary fines on peasants wherever land was sold or inherited. However over the $15^{th}$ and $16^{th}$ centuries landlords  successfully crushed peasant land rights in order to combat this decreasing power. This decay in freehold rights meant that lord could slowly accumulate large demesne holdings, which could be leased out to larger tenants \cite[p.7]{hoyle1990tenure}. Who in turn sub-let to smaller farmers or employed wage labourers. This new wave of larger tenancies brought two new sources of market imperatives, firstly the lord had to keep his rents competitive as he was now trying to undercut the rest of the domestic tenancy market, given that the larger tenant had the ability to move away. Secondly the tenant had to compete on a goods market, requiring them to sell the products of their husbandry in order to make enough profit to fulfil market dependant rent. This required a shift from production for subsistence to production for sale. These two forces resulted in a new weakly symbiotic relation between lord and capitalist tenant in which both were driven to re-invest surplus value back into their land, thereby improving agricultural yields and transforming feudal wealth into capital. This established a new set of laws of motion in which the lord and tenant were compelled to improve their agricultural yields in order to self-replicate their class condition. Brenner claims that ``agricultural improvement not only made it possible for an ever greater proportion of the population to leave the land to enter industry" \cite[p.67]{brenner1976agrarian} but also mean that a larger proportion of the domestic population was no longer involved in agriculture, freeing peasants to become the wage labour required for ``England's continued industrial growth" \cite[p.67]{brenner1976agrarian} through the $17^{th}$ Century. 
\section{Model Design}
\section{Results}
\section{Discussion}
\section{Conclusions}

\section{Acknowledgements}
This work has received funding from the European Union’s Horizon 2020 research and
innovation programme under the Marie Skłodowska-Curie grant agreement No
956107, ``Economic Policy in Complex Environments (EPOC)”.

\bibliographystyle{apalike}
\bibliography{bib}
\end{multicols}
\end{document}
